% Mouvement de chute libre d'une bille dans un plan
On lance, de nouveau, à partir de la position $O$, la bille $(S)$ (précédente)
avec une vitesse initiale $\vec{v}_{02}$ faisant un angle $\alpha$ avec
l'horizontale. On étudie le mouvement de $G$, centre d'inertie de la bille $S$,
dans un repère orthonormé $(O,\vec{i},\vec{j})$ lié à la Terre, considéré
comme galiléen (figure 2).

\subsection*{1.2}
En appliquant la deuxième loi de Newton, trouver l'expression littérale des équations horaires $x(t)$ et $y(t)$ du mouvement de $G$.

\subsection*{2.2}
Montrer que l'expression de la portée est :
\[
x_p = \frac{v_{02}^2 \cdot \sin(2\alpha)}{g}
\]

\subsection*{3.2}
À l'aide d'un logiciel approprié, on a obtenu le document de la figure (3) représentant les trajectoires du mouvement de $G$ pour la même valeur de la vitesse initiale $v_{02}$ et pour différents angles de tir $\alpha_0 = 45^\circ$, $\alpha_1$ et $\alpha_2$.

\begin{center}
\textbf{Figure 3}
\end{center}

\begin{enumerate}
    \item[\textbf{a.}] Déterminer la valeur de la portée $x_p$ correspondant à l'angle de tir $\alpha_0$.
    \item[\textbf{b.}] Déterminer la valeur de l'angle $\alpha_1$. En déduire la valeur de l'angle $\alpha_2$.
    \item[\textbf{c.}] Lorsque la vitesse de $G$ est $v_1$ pour l'angle de tir $\alpha_1$ et $v_2$ pour l'angle de tir $\alpha_2$, laquelle des relations suivantes est correcte ? Recopier la bonne réponse sur votre feuille.
\end{enumerate}

\begin{center}
\textbf{Figure 3 (suite)}
\end{center}

\begin{itemize}
    \item[\textbf{a.}] $v_1 = 3,2 \, v_2$
    \item[\textbf{b.}] $v_1 = 1,6 \, v_2$
    \item[\textbf{c.}] $v_1 = 0,8 \, v_2$
    \item[\textbf{d.}] $v_1 = 0,4 \, v_2$
\end{itemize}

